% =====================================================================
%  Hp-DAC — LaTeX TABLE STANDARD
%
%  Self-contained. Upload to Overleaf and press Recompile: the PDF shows
%  every template rendered, and the source below each banner is what you
%  copy into your own document.
%
%  To use in a real manuscript:
%    1. copy the PREAMBLE block into your own preamble;
%    2. copy the template closest to what you need;
%    3. edit the caption, label, column format and data.
%
%  Project: Hp-DAC (ANP / Repsol Sinopec Brasil, POLO-UFSC)
% =====================================================================

\documentclass[11pt,a4paper]{article}

% ---------------------------------------------------------------------
%  PREAMBLE — copy this block into your manuscript
% ---------------------------------------------------------------------
\usepackage[T1]{fontenc}
\usepackage[utf8]{inputenc}
\usepackage{lmodern}          % same font family as the figures
\usepackage{amsmath}

\usepackage{booktabs}         % \toprule \midrule \bottomrule \cmidrule
\usepackage{array}            % column type tuning
\usepackage{siunitx}          % number alignment and units
\usepackage{threeparttable}   % table notes tied to the table
\usepackage{longtable}        % tables that break across pages
\usepackage{pdflscape}        % landscape page for very wide tables
\usepackage{caption}

%  per-mode=power is the important one: it typesets \unit{\watt\per\metre}
%  as W m^{-1}, matching the negative exponents used on the figure axes.
%  A solidus is banned in this project because it is ambiguous with two
%  denominators: "W/m K" can be read two ways.
\sisetup{
  per-mode            = power,
  exponent-product    = \times,
  separate-uncertainty= true,   % 5.27 +- 0.15 rather than 5.27(15)
  detect-all,                   % inherit the surrounding font
  table-align-uncertainty = true,
  group-digits        = integer,
  group-minimum-digits= 5,      % 1234 stays 1234; 12345 becomes 12 345
}

\captionsetup[table]{skip=6pt}  % caption sits above the table
% ---------------------------------------------------------------------

\usepackage[margin=25mm]{geometry}
\usepackage{hyperref}

\title{Hp-DAC --- LaTeX Table Standard}
\date{}

\begin{document}
\maketitle

\section*{House rules}

\begin{enumerate}
  \item \textbf{\texttt{booktabs} rules only.} One \verb|\toprule|, one
        \verb|\midrule| under the header, one \verb|\bottomrule|.
        \emph{No vertical rules, ever. No \texttt{\textbackslash hline}.}
        A vertical rule almost always means the table should be split or
        transposed.
  \item \textbf{Caption above the table}, sentence case, ending in a full
        stop. Figures take the caption below; tables above.
  \item \textbf{Header carries the symbol and the unit} in square brackets,
        dimensionless as [--]. Units use negative exponents:
        \verb|[kg h$^{-1}$]|, never \verb|[kg/h]|. Keep the symbol identical
        to the one used in the figures and the nomenclature list.
  \item \textbf{Decimals match the measurement} and stay consistent down a
        column: \num{5.27} and \num{4.32}, never \num{5.27} and 4.3.
        An uncertainty carries the same decimals as its value.
  \item \textbf{Missing entries are \texttt{--}}, never blank.
  \item \textbf{Label as \texttt{tab:name}}, hyphenated, matching the file
        name if the table lives in its own file.
  \item \textbf{Never use \texttt{\textbackslash resizebox} on a table.}
        Shrinking the font to fit is a sign of too many columns: split the
        table, transpose it, or use the landscape template.
\end{enumerate}

\clearpage
% =====================================================================
%  TEMPLATE 1 — BASIC RESULTS TABLE
%  The default. Use it unless you need one of the others.
%  Column format: l for text, c for numbers of equal width.
% =====================================================================
\section*{Template 1 --- basic results table}

\begin{table}[!htb]
  \centering
  \caption{Mean performance of each prototype version.}
  \label{tab:mean-values}
  \begin{tabular}{lccc}
    \toprule
    Version & $\dot{m}_\mathrm{f}$ [kg h$^{-1}$]
            & $COP$ [--]
            & $\eta_\mathrm{II}$ [--] \\
    \midrule
    1 &  834.2 & 5.27 & 0.50 \\
    2 &  845.1 & 4.32 & 0.60 \\
    3 & 1219.4 & 2.68 & 0.70 \\
    4 & 1326.0 & 1.59 & 0.80 \\
    5 & 1541.3 & 0.79 & 0.90 \\
    \bottomrule
  \end{tabular}
\end{table}

% =====================================================================
%  TEMPLATE 2 — DECIMAL ALIGNMENT WITH siunitx
%  Preferred whenever a column holds numbers with different digit counts.
%
%  S[table-format=4.1] reserves 4 integer digits and 1 decimal.
%  CATCH: a header inside an S column MUST be wrapped in braces {...},
%  otherwise siunitx tries to parse it as a number and the compile fails
%  with "Missing number, treated as zero". This is the most common error.
% =====================================================================
\section*{Template 2 --- decimal alignment}

\begin{table}[!htb]
  \centering
  \caption{Heating capacity and coefficient of performance, decimal
           aligned on the S columns.}
  \label{tab:aligned}
  \begin{tabular}{l
                  S[table-format=4.1]
                  S[table-format=1.2]
                  S[table-format=2.2]}
    \toprule
    Version & {$\dot{m}_\mathrm{f}$ [kg h$^{-1}$]}
            & {$COP$ [--]}
            & {$\dot{Q}_\mathrm{H}$ [kW]} \\
    \midrule
    1 &  834.2 & 5.27 & 12.40 \\
    2 &  845.1 & 4.32 & 18.03 \\
    3 & 1219.4 & 2.68 & 24.60 \\
    4 & 1326.0 & 1.59 & 31.17 \\
    5 & 1541.3 & 0.79 & 38.90 \\
    \bottomrule
  \end{tabular}
\end{table}

\clearpage
% =====================================================================
%  TEMPLATE 3 — MEASUREMENTS WITH UNCERTAINTY
%  table-format=1.2(2) means 1 integer digit, 2 decimals, 2 uncertainty
%  digits. With separate-uncertainty=true (set in the preamble) siunitx
%  prints 5.27 +- 0.15 and aligns the +- signs.
%
%  Write the value in the source as 5.27(15) -- the parenthesis carries
%  the uncertainty in units of the last decimal.
% =====================================================================
\section*{Template 3 --- uncertainties}

\begin{table}[!htb]
  \centering
  \caption{Measured coefficient of performance with expanded uncertainty
           ($k = 2$).}
  \label{tab:uncertainty}
  \begin{tabular}{l
                  S[table-format=4.0]
                  S[table-format=1.2(2)]
                  S[table-format=2.1]}
    \toprule
    Version & {$\dot{m}_\mathrm{f}$ (kg h$^{-1}$)}
            & {$COP$ (--)}
            & {$\Delta T_\mathrm{lift}$ (K)} \\
    \midrule
    1 &  834 & 5.27(15) & 65.0 \\
    2 &  845 & 4.32(12) & 68.4 \\
    3 & 1219 & 2.68(9)  & 71.2 \\
    4 & 1326 & 1.59(7)  & 74.8 \\
    5 & 1541 & 0.79(5)  & 78.1 \\
    \bottomrule
  \end{tabular}
\end{table}

% =====================================================================
%  TEMPLATE 4 — GROUPED COLUMN HEADERS
%  \cmidrule(lr){2-3} draws a partial rule under a group. The (lr) trims
%  it so adjacent groups do not touch. Never use a full-width \midrule
%  inside the body to separate groups.
%
%  NOTE the last column: table-format=+2.1 reserves a slot for the minus
%  sign. Without the +, a negative value is wider than declared and LaTeX
%  reports "Overfull \hbox" -- the usual cause of that warning in a table.
% =====================================================================
\section*{Template 4 --- grouped headers}

\begin{table}[!htb]
  \centering
  \caption{Measured and predicted performance for each prototype version.}
  \label{tab:grouped}
  \begin{tabular}{l
                  S[table-format=1.2] S[table-format=2.2]
                  S[table-format=1.2] S[table-format=2.2]
                  S[table-format=+2.1]}
    \toprule
    & \multicolumn{2}{c}{Measured} & \multicolumn{2}{c}{Predicted}
    & {Deviation} \\
    \cmidrule(lr){2-3} \cmidrule(lr){4-5}
    Version & {$COP$ (--)} & {$\dot{Q}_\mathrm{H}$ (kW)}
            & {$COP$ (--)} & {$\dot{Q}_\mathrm{H}$ (kW)}
            & {(\%)} \\
    \midrule
    1 & 5.27 & 12.40 & 5.19 & 12.61 &  1.7 \\
    2 & 4.32 & 18.03 & 4.41 & 17.88 & -2.1 \\
    3 & 2.68 & 24.60 & 2.63 & 24.92 &  1.9 \\
    4 & 1.59 & 31.17 & 1.64 & 30.85 & -3.1 \\
    5 & 0.79 & 38.90 & 0.77 & 39.44 &  2.5 \\
    \bottomrule
  \end{tabular}
\end{table}

\clearpage
% =====================================================================
%  TEMPLATE 5 — PARAMETERS TABLE WITH NOTES
%  For model inputs, boundary conditions and assumptions. Mixed text and
%  numbers, so plain c columns rather than S.
%
%  threeparttable keeps the notes the same width as the table and ties
%  \tnote{a} in the body to \item[a] in the notes.
% =====================================================================
\section*{Template 5 --- parameters, with table notes}

\begin{table}[!htb]
  \centering
  \begin{threeparttable}
    \caption{Boundary conditions and model parameters used in the
             simulations.}
    \label{tab:parameters}
    \begin{tabular}{llcc}
      \toprule
      Parameter & Symbol & Value & Unit \\
      \midrule
      Heat source inlet temperature   & $T_\mathrm{source}$ & 25.0 & $^{\circ}$C \\
      Heat sink outlet temperature    & $T_\mathrm{sink}$   & 90.0 & $^{\circ}$C \\
      Isentropic efficiency\tnote{a}  & $\eta_\mathrm{is}$  & 0.70 & -- \\
      Capture capacity                & $\dot{m}_\mathrm{CO_2}$ & 5000 & t yr$^{-1}$ \\
      Discount rate\tnote{b}          & $r$                 & 0.08 & -- \\
      Plant lifetime                  & $n$                 & 20   & yr \\
      Pressure drop                   & $\Delta p$          & --   & kPa \\
      \bottomrule
    \end{tabular}
    \begin{tablenotes}[flushleft]\footnotesize
      \item[a] Assumed constant over the operating envelope.
      \item[b] Taken from the optimised case of Teixeira and Nunes (2024).
    \end{tablenotes}
  \end{threeparttable}
\end{table}

% =====================================================================
%  TEMPLATE 6 — TABLE THAT BREAKS ACROSS PAGES
%  longtable is its own environment: no surrounding \begin{table}, and
%  the caption goes inside. \endfirsthead / \endhead repeat the header.
% =====================================================================
\section*{Template 6 --- multi-page table}

\begin{longtable}{l S[table-format=4.1] S[table-format=1.2]}
  \caption{Full test matrix.} \label{tab:long} \\
  \toprule
  Test & {$\dot{m}_\mathrm{f}$ (kg h$^{-1}$)} & {$COP$ (--)} \\
  \midrule
  \endfirsthead
  \multicolumn{3}{l}{\footnotesize\itshape Table \ref{tab:long}, continued}\\
  \toprule
  Test & {$\dot{m}_\mathrm{f}$ (kg h$^{-1}$)} & {$COP$ (--)} \\
  \midrule
  \endhead
  \midrule
  \multicolumn{3}{r}{\footnotesize\itshape continued on the next page}\\
  \endfoot
  \bottomrule
  \endlastfoot
  T01 &  834.2 & 5.27 \\
  T02 &  845.1 & 4.32 \\
  T03 & 1219.4 & 2.68 \\
  T04 & 1326.0 & 1.59 \\
  T05 & 1541.3 & 0.79 \\
\end{longtable}

% =====================================================================
%  TEMPLATE 7 — VERY WIDE TABLE
%  Rotate the page rather than shrinking the font. pdflscape rotates the
%  PDF page itself, so the table is readable on screen without tilting
%  your head. Use this instead of \resizebox.
%
%  Count the digits before writing table-format: T_sink reaches 120.0, so
%  it is 3.1 and not 2.1. Declaring too few digits is the second common
%  cause of an "Overfull \hbox" warning.
% =====================================================================
\begin{landscape}
\section*{Template 7 --- wide table on a landscape page}

\begin{table}[!htb]
  \centering
  \caption{Full performance map of the heat pump across the operating
           envelope.}
  \label{tab:wide}
  \begin{tabular}{l
                  S[table-format=2.1] S[table-format=3.1]
                  S[table-format=1.2] S[table-format=2.2]
                  S[table-format=2.2] S[table-format=1.3]
                  S[table-format=3.1]}
    \toprule
    Case & {$T_\mathrm{source}$ ($^{\circ}$C)}
         & {$T_\mathrm{sink}$ ($^{\circ}$C)}
         & {$COP$ (--)}
         & {$\dot{Q}_\mathrm{H}$ (kW)}
         & {$\dot{W}_\mathrm{cp}$ (kW)}
         & {$\eta_\mathrm{II}$ (--)}
         & {$\Delta T_\mathrm{lift}$ (K)} \\
    \midrule
    A & 25.0 & 80.0 & 3.42 & 12.40 &  3.63 & 0.512 &  55.0 \\
    B & 25.0 & 90.0 & 2.87 & 18.03 &  6.28 & 0.478 &  65.0 \\
    C & 35.0 & 90.0 & 3.16 & 24.60 &  7.78 & 0.495 &  55.0 \\
    D & 35.0 & 110.0 & 2.24 & 31.17 & 13.92 & 0.441 &  75.0 \\
    E & 45.0 & 120.0 & 2.05 & 38.90 & 18.98 & 0.428 &  75.0 \\
    \bottomrule
  \end{tabular}
\end{table}
\end{landscape}

\clearpage
\section*{Quick reference}

\begin{table}[!htb]
  \centering
  \caption{Column types and when to use them.}
  \label{tab:column-types}
  \begin{tabular}{ll}
    \toprule
    Column format & Use for \\
    \midrule
    \verb|l|                        & Text: case names, versions, descriptions \\
    \verb|c|                        & Numbers already of equal width \\
    \verb|S[table-format=4.1]|      & Numbers to align on the decimal point \\
    \verb|S[table-format=1.2(2)]|   & Value with uncertainty \\
    \verb|S[table-format=2.1e1]|    & Scientific notation \\
    \verb|p{30mm}|                  & Wrapped text of fixed width \\
    \bottomrule
  \end{tabular}
\end{table}

\noindent
\textbf{Do not use:} \verb!|! between columns, \verb|\hline|,
\verb|\resizebox|, \verb|\small| to make a table fit, or a caption below
the table.

\bigskip
\noindent
\textbf{Units, quick check.} Write \verb|(kg h$^{-1}$)| in a plain header,
or \verb|\unit{\kilogram\per\hour}| if you prefer siunitx to typeset it ---
with \verb|per-mode=power| both give the same result. Never write
\verb|(kg/h)|.

\end{document}
